\documentclass[a4paper, 12pt]{article}
\usepackage{deuresutils}

\title{Cosararar}
\asignatura{Mètodes Numèrics i Probabilistics}
\author{Eduardo Pérez Motato}
\niu{NIU: 1709992}
\date{28/03/2025}

\begin{document}
    \makeheader

    \begin{exercici}
    	Considereu el mètode d'integració numèrica donat per la formula
        \begin{displaymath}
        	I(f) = \int_{-1}^{1} f(x) \, dx \approx S(f) = \alpha f(-1) + \beta f(1) + \gamma f'(-1) + \delta f'(1)
        \end{displaymath}
        \begin{enumerate}
        	\item Determinar els pesos de $\alpha, \beta, \gamma, \delta$ de l'aproximació $S(f)$ de manera que sigui exacta per a polinomis de grau el més gran possible. Quin és el grau màxim?
         	\item Comproveu per $f(x) = \frac{1}{3+x}$ i comproveu amb el valor exacte.
        \end{enumerate}
    \end{exercici}
    \begin{solucio}
    	Per a determinar els pesos $\alpha, \beta, \gamma, \delta$ de l'aproximació $S(f)$, fem servir $f(x)$ que sigui igual a $1, x, x^2, x^3$.
     	Llavors tenim el següent sistema:
      	\begin{displaymath}
	      	\begin{cases}
	       	    \alpha + \beta + \gamma + \delta = 2 \\
				-\alpha + \beta - \gamma + \delta = 0 \\
				\alpha + \beta - \gamma - \delta = \frac{2}{3} \\
				-\alpha + \beta + \gamma - \delta = 0
	       	\end{cases}
	  	\end{displaymath}
        Resolent aquest sistema obtenim $\alpha = \beta = 1$ i $\gamma = -\delta = \frac{1}{3}$.
		Per això com a mínim serà exacte per a polinomis de grau 3. Comprovem per $f(x) = x^4$ si es compleix:
		\begin{displaymath}
			\int_{-1}^{1} x^4 \, dx \neq \alpha+\beta-4\gamma+4\gamma = 2
		\end{displaymath}
		Com no es compleix, és exacte com a màxim per cúbic.\\
		Per comprovar amb $f(x) = \frac{1}{3+x}$, tenim:
		\begin{displaymath}
			\int_{-1}^{1} \frac{1}{3+x} \, dx = \ln(2) = S\left(\frac{1}{3+x}\right) = \frac{1}{3} - \frac{1}{3} + \frac{1}{6} - \frac{1}{6} = 0
		\end{displaymath}
		Per tant, no es compleix.
    \end{solucio}
\end{document}
